%%%%%%%%%%%%%%%%%%%%%%%%%%%%%%%%%%%%%%%%%%%%%%%%%%%%%%%%%%%%%%%%%%%%%%%%%%%%
% Exploratory Pilot Evaluation -- SafeBARS ethical mirror "finds" ethical
% blind spots. Structured after Taheri et al. (CHI'26), "I followed what
% felt right, not what I was told": Autonomy, Coaching, and Recognizing
% Bias Through AI-Mediated Dialogue.
% HONESTY NOTE: n=2 pilot, real session archives only. No fabricated n.
%%%%%%%%%%%%%%%%%%%%%%%%%%%%%%%%%%%%%%%%%%%%%%%%%%%%%%%%%%%%%%%%%%%%%%%%%%%%

\subsection{Exploratory Pilot Evaluation: Does the Mirror Surface Blind Spots Researchers Miss?}
\label{sec:pilot}

Before scaling to a controlled study (\S\ref{sec:evaldesign}), we ran two
exploratory pilot sessions to test the core mechanism claim: that the
mirror surfaces ethical problems a researcher did \emph{not} include in their
own written plan. This parallels the effectiveness question in Taheri et
al.~\cite{taheri26} (whether AI-mediated dialogue improves recognition of
bias), but where they measured recognition via pre/post vignette ratings, we
measure \emph{discovery yield}---the count of distinct ethical items the
mirror surfaces, the blind-spot rate, and the coverage gap against the
researcher's plan---because at pilot scale ($n=2$) inference statistics are
not meaningful. All counts and quotations below are drawn verbatim from the
two real session archives.\footnote{Archives: \texttt{SafeBARS\_StudentTest/study\_bundle.json}
(R1, full guided session) and \texttt{SafeBARS\_StudentTest\_2} (R2, Deep-audit
verification run).}

\paragraph{Method.}
Two researchers (R1, R2) entered the mirror with their own un-reviewed
research concepts: R1, \emph{fairness in shared dorm living} (an AI that
reads smart-plug electricity and group-chat tone to label a ``free-riding''
roommate and draft a private message); R2, \emph{campus mental-health
crisis detection (CampusMind)} (an AI that scans anonymized forum/group-chat
posts to quietly connect struggling students to a counseling center). Both
followed the guided protocol (lens selection $\rightarrow$ role probes
$\rightarrow$ stakeholder map $\rightarrow$ dissonance edges $\rightarrow$
revision ledger); R2's session additionally ran the Deep-audit stress test.
Materials were the nine-lens ethical framework and the Deep-audit module
(internal-contradiction detection, additional dissonance edges, real-world
failure analogues, forced-reflection lenses).

\paragraph{5.1 Discovery yield (quantitative).}
Table~\ref{tab:pilot} summarizes the real discovery counts. In both sessions
the mirror surfaced \textbf{$\ge 9$ distinct ethical items}; of these,
\textbf{$\ge 3$ were blind spots} absent from or merely declared (not
action-linked) in the original plan. In R1's session, \textbf{6/9 lenses
were not action-linked} in the plan (Figure~\ref{fig:pilotcov})---the mirror
exposed gaps the researcher had left open. R2's Deep-audit additionally
surfaced \textbf{1 internal contradiction + 3 extra dissonance edges + 5
real-world failure analogues}---dimensions entirely absent from the written
plan. R1 engaged with every tension the mirror raised (2/2 decisions,
100\%).

\begin{table}[t]
  \centering
  \small
  \caption{Pilot discovery yield (real counts from two archived sessions).}
  \label{tab:pilot}
  \begin{tabular}{p{1.5cm}p{2.6cm}p{1.4cm}p{2.4cm}p{1.1cm}}
    \toprule
    Session & Research domain & Lenses surfaced & Blind spots vs. plan & Tension edges \\
    \midrule
    R1 Dorm & Shared-dorm fairness & 9 (2/4/2/1\textsuperscript{a}) & $\ge$3 (downstream misuse, redress, analogues) & 2 \\
    R2 CampusMind & Crisis detection & 6 forced + 9 & $\ge$3 (no appeal, reuse cap, consent contradiction) & 4 \\
    \bottomrule
    \multicolumn{5}{l}{\footnotesize\textsuperscript{a} coverage states missing/claimed/reasoned/action-linked.}
  \end{tabular}
\end{table}

\begin{figure}[t]
  \centering
  \includegraphics[width=0.46\columnwidth]{fig_pilot_coverage.pdf}
  \caption{Dorm pilot: coverage states of the 9 lenses the mirror derived from
  the researcher's own plan. Red = missing, orange = claimed-only, blue =
  reasoned, green = action-linked. 6/9 were not action-linked in the
  original plan---i.e. the mirror surfaced ethical gaps the researcher had
  not closed.}
  \label{fig:pilotcov}
\end{figure}

\paragraph{5.2 Qualitative findings.}
Reflexive thematic analysis of the two sessions yielded five themes;
quotations are verbatim from the archives.

\noindent\textbf{Theme A -- Downstream misuse \& scale spillover (plan gap).}
The mirror flagged that R1's plan ignored repurposing of its outputs:
\textit{``the plan does not consider how outputs (e.g., drafted messages,
inferred `free-rider' labels) could be repurposed outside the intended
context---e.g., shared with landlords, cited in conduct hearings, or
integrated into university disciplinary systems''}
(lens \texttt{downstream\_use\_misuse\_scale}, status = missing).

\noindent\textbf{Theme B -- Contestability / appeal path missing.}
R2's Deep-audit marked AUD-EDGE-003 (severity = elevated): \textit{``The plan
monitors or flags people but states no appeal or redress path.''} and
AUD-EDGE-002 (severity = high): \textit{``The system is built for one caring
purpose but has no stated limit preventing administrative or disciplinary
reuse.''} Neither safeguard appeared in the researcher's plan.

\noindent\textbf{Theme C -- Internal contradiction the researcher missed.}
R2's CNTR-001 exposed a logical conflict between the plan's promise of
opt-out non-disclosure and its silent-notification behavior---a contradiction
the original tool (which only reflects what is written) would have missed,
and which Deep-audit's added capability caught.

\noindent\textbf{Theme D -- Real-world failure analogues as evidence bridges.}
The mirror automatically bridged five real failures, turning abstract worry
into citable precedent, e.g. \textit{exam\_proctoring\_surveillance\_2020}:
\textit{``Surveillance tooling aimed at `integrity' can punish the already
marginalised. Offer a human-reviewed alternative and disclose the
monitoring.''}; \textit{obermeyer\_health\_bias\_2019}: \textit{``A proxy
variable that looks neutral can silently encode the very bias you are trying
to avoid. Audit subgroup disparity, not just overall accuracy.''};
\textit{uk\_a\_levels\_algorithm\_2020}: \textit{``When an algorithm
overrides human assessment at scale, small model assumptions produce large,
visible injustices and loss of public trust.''}

\noindent\textbf{Theme E -- Researcher engagement with findings.}
R1 acted on both tension edges: \texttt{EDGE-001} $\rightarrow$
\texttt{add\_safeguard} (``Assign a specific role for post-update conflict
auditing'') and \texttt{EDGE-002} $\rightarrow$ \texttt{consult\_stakeholders}.
Decision rate 100\%, indicating the surfaced problems were received and
responded to rather than ignored.

\paragraph{What this pilot supports, and what it does not.}
The pilot \textbf{supports} the claim that \emph{SafeBARS surfaces ethical
problems researchers did not include or did not operationalize in their own
plans}, evidenced by real per-session counts ($\ge$9 items, $\ge$3 blind
spots, an internal contradiction, and five analogues) and verbatim
quotations---structurally parallel to the qualitative findings in Taheri et
al. A verbatim replay of R1's mirror conversation (Appendix~\ref{fig:replay})
shows the AI surfacing each tension step by step from the researcher's own
plan. It \textbf{does not yet support} a claim that the mirror \emph{corrects}
behavior: R1's ledger recorded decisions but noted \textit{``No concise
changed passage was recorded''}, i.e. the change to the plan text was not
captured. That gap is exactly what the scaled Study~1 (\S\ref{sec:evaldesign})
is designed to measure, by recording before/after plan diffs. Finally, with
$n=2$ this pilot cannot support inferential conclusions; the controlled
Study~1 (multimodal mirror vs.\ text-only baseline, $N=80$--$120$) will
provide Taheri-style quantitative evidence (discovery-yield difference,
blind-spot detection rate, effect sizes).
